% Converted from S5_Tool_call_audit/reviewer_response_package_20261119_claude_v13/figures/FIGURE_CAPTION_EN.md and S5_Tool_call_audit/reviewer_response_package_20261119_claude_v13/figures/Fig_R2_16_traces.pdf (2026-11-17)
\begin{responsefigure}
\includegraphics[width=\textwidth,height=0.64\textheight,keepaspectratio]{report_rep_execution_traces.pdf}
\caption{\responsefiguretitle{fig_rep_execution_traces}{Representative execution traces of the planner--executor workflow. Each row shows one run from the audit (150 cases $\times$ 3 runs). Boxes are tool calls in the recorded execution order, given by the small numbers; arrows indicate order only, not data flow. The text under each box summarises the recorded arguments and returns, and the left column gives the dataset, the number of input images and the observed outcome. Fill colour and symbol together encode the execution status (legend). \textbf{a}, No thyroid nodule. No image was flagged by preprocessing, so segmentation, feature classification and fusion were not scheduled. Lymph-node screening of the four neck images returned one positive image, which the report reflects. \textbf{b}, Multiple nodules. Two preprocessing-flagged images were confirmed by segmentation, classified, and fused into two nodules, each reported with a size. \textbf{c}, Conflicting intermediate evidence. Of three preprocessing-flagged images, two were confirmed by segmentation; the third was not passed to classification or fusion. The same disagreement occurred in all three runs of this case and of one other case. \textbf{d}, Invalid tool name and recovery. The LLM emitted a malformed tool name and the executor returned an unknown-tool error; the next call used the registered name with identical arguments. \textbf{e}, Off-region targets not processed. A lymph-node screening request listed five images whose preprocessing region was not the neck. The executor passed only the neck images to the screening tool; the other five were not processed in this call. \textbf{f}, Dependency failure. Fusion ran before pixel-spacing estimates were available, so the fused nodules had no millimetre sizes. Spacing was computed later, but fusion was not repeated, and the report contains two nodule descriptions without sizes. \textbf{g}, Omission. A nodule was confirmed by segmentation and classified, but fusion was neither planned nor executed (dashed box), and the report contains no nodule description. Plan-guided recovery could not add the step because it was absent from the plan. Frequencies of each failure mode are given in \responsetableref{tab_rep_failure_modes}. Case identifiers, file paths and report text are omitted. LN, lymph node.}}
\label{fig_rep_execution_traces}
\end{responsefigure}
